\documentclass[11pt]{article}

    \usepackage[breakable]{tcolorbox}
    \usepackage{parskip} % Stop auto-indenting (to mimic markdown behaviour)
    

    % Basic figure setup, for now with no caption control since it's done
    % automatically by Pandoc (which extracts ![](path) syntax from Markdown).
    \usepackage{graphicx}
    % Keep aspect ratio if custom image width or height is specified
    \setkeys{Gin}{keepaspectratio}
    % Maintain compatibility with old templates. Remove in nbconvert 6.0
    \let\Oldincludegraphics\includegraphics
    % Ensure that by default, figures have no caption (until we provide a
    % proper Figure object with a Caption API and a way to capture that
    % in the conversion process - todo).
    \usepackage{caption}
    \DeclareCaptionFormat{nocaption}{}
    \captionsetup{format=nocaption,aboveskip=0pt,belowskip=0pt}

    \usepackage{float}
    \floatplacement{figure}{H} % forces figures to be placed at the correct location
    \usepackage{xcolor} % Allow colors to be defined
    \usepackage{enumerate} % Needed for markdown enumerations to work
    \usepackage{geometry} % Used to adjust the document margins
    \usepackage{amsmath} % Equations
    \usepackage{amssymb} % Equations
    \usepackage{textcomp} % defines textquotesingle
    % Hack from http://tex.stackexchange.com/a/47451/13684:
    \AtBeginDocument{%
        \def\PYZsq{\textquotesingle}% Upright quotes in Pygmentized code
    }
    \usepackage{upquote} % Upright quotes for verbatim code
    \usepackage{eurosym} % defines \euro

    \usepackage{iftex}
    \ifPDFTeX
        \usepackage[T1]{fontenc}
        \IfFileExists{alphabeta.sty}{
              \usepackage{alphabeta}
          }{
              \usepackage[mathletters]{ucs}
              \usepackage[utf8x]{inputenc}
          }
    \else
        \usepackage{fontspec}
        \usepackage{unicode-math}
    \fi

    \usepackage{fancyvrb} % verbatim replacement that allows latex
    \usepackage{grffile} % extends the file name processing of package graphics
                         % to support a larger range
    \makeatletter % fix for old versions of grffile with XeLaTeX
    \@ifpackagelater{grffile}{2019/11/01}
    {
      % Do nothing on new versions
    }
    {
      \def\Gread@@xetex#1{%
        \IfFileExists{"\Gin@base".bb}%
        {\Gread@eps{\Gin@base.bb}}%
        {\Gread@@xetex@aux#1}%
      }
    }
    \makeatother
    \usepackage[Export]{adjustbox} % Used to constrain images to a maximum size
    \adjustboxset{max size={0.9\linewidth}{0.9\paperheight}}

    % The hyperref package gives us a pdf with properly built
    % internal navigation ('pdf bookmarks' for the table of contents,
    % internal cross-reference links, web links for URLs, etc.)
    \usepackage{hyperref}
    % The default LaTeX title has an obnoxious amount of whitespace. By default,
    % titling removes some of it. It also provides customization options.
    \usepackage{titling}
    \usepackage{longtable} % longtable support required by pandoc >1.10
    \usepackage{booktabs}  % table support for pandoc > 1.12.2
    \usepackage{array}     % table support for pandoc >= 2.11.3
    \usepackage{calc}      % table minipage width calculation for pandoc >= 2.11.1
    \usepackage[inline]{enumitem} % IRkernel/repr support (it uses the enumerate* environment)
    \usepackage[normalem]{ulem} % ulem is needed to support strikethroughs (\sout)
                                % normalem makes italics be italics, not underlines
    \usepackage{soul}      % strikethrough (\st) support for pandoc >= 3.0.0
    \usepackage{mathrsfs}
    

    
    % Colors for the hyperref package
    \definecolor{urlcolor}{rgb}{0,.145,.698}
    \definecolor{linkcolor}{rgb}{.71,0.21,0.01}
    \definecolor{citecolor}{rgb}{.12,.54,.11}

    % ANSI colors
    \definecolor{ansi-black}{HTML}{3E424D}
    \definecolor{ansi-black-intense}{HTML}{282C36}
    \definecolor{ansi-red}{HTML}{E75C58}
    \definecolor{ansi-red-intense}{HTML}{B22B31}
    \definecolor{ansi-green}{HTML}{00A250}
    \definecolor{ansi-green-intense}{HTML}{007427}
    \definecolor{ansi-yellow}{HTML}{DDB62B}
    \definecolor{ansi-yellow-intense}{HTML}{B27D12}
    \definecolor{ansi-blue}{HTML}{208FFB}
    \definecolor{ansi-blue-intense}{HTML}{0065CA}
    \definecolor{ansi-magenta}{HTML}{D160C4}
    \definecolor{ansi-magenta-intense}{HTML}{A03196}
    \definecolor{ansi-cyan}{HTML}{60C6C8}
    \definecolor{ansi-cyan-intense}{HTML}{258F8F}
    \definecolor{ansi-white}{HTML}{C5C1B4}
    \definecolor{ansi-white-intense}{HTML}{A1A6B2}
    \definecolor{ansi-default-inverse-fg}{HTML}{FFFFFF}
    \definecolor{ansi-default-inverse-bg}{HTML}{000000}

    % common color for the border for error outputs.
    \definecolor{outerrorbackground}{HTML}{FFDFDF}

    % commands and environments needed by pandoc snippets
    % extracted from the output of `pandoc -s`
    \providecommand{\tightlist}{%
      \setlength{\itemsep}{0pt}\setlength{\parskip}{0pt}}
    \DefineVerbatimEnvironment{Highlighting}{Verbatim}{commandchars=\\\{\}}
    % Add ',fontsize=\small' for more characters per line
    \newenvironment{Shaded}{}{}
    \newcommand{\KeywordTok}[1]{\textcolor[rgb]{0.00,0.44,0.13}{\textbf{{#1}}}}
    \newcommand{\DataTypeTok}[1]{\textcolor[rgb]{0.56,0.13,0.00}{{#1}}}
    \newcommand{\DecValTok}[1]{\textcolor[rgb]{0.25,0.63,0.44}{{#1}}}
    \newcommand{\BaseNTok}[1]{\textcolor[rgb]{0.25,0.63,0.44}{{#1}}}
    \newcommand{\FloatTok}[1]{\textcolor[rgb]{0.25,0.63,0.44}{{#1}}}
    \newcommand{\CharTok}[1]{\textcolor[rgb]{0.25,0.44,0.63}{{#1}}}
    \newcommand{\StringTok}[1]{\textcolor[rgb]{0.25,0.44,0.63}{{#1}}}
    \newcommand{\CommentTok}[1]{\textcolor[rgb]{0.38,0.63,0.69}{\textit{{#1}}}}
    \newcommand{\OtherTok}[1]{\textcolor[rgb]{0.00,0.44,0.13}{{#1}}}
    \newcommand{\AlertTok}[1]{\textcolor[rgb]{1.00,0.00,0.00}{\textbf{{#1}}}}
    \newcommand{\FunctionTok}[1]{\textcolor[rgb]{0.02,0.16,0.49}{{#1}}}
    \newcommand{\RegionMarkerTok}[1]{{#1}}
    \newcommand{\ErrorTok}[1]{\textcolor[rgb]{1.00,0.00,0.00}{\textbf{{#1}}}}
    \newcommand{\NormalTok}[1]{{#1}}

    % Additional commands for more recent versions of Pandoc
    \newcommand{\ConstantTok}[1]{\textcolor[rgb]{0.53,0.00,0.00}{{#1}}}
    \newcommand{\SpecialCharTok}[1]{\textcolor[rgb]{0.25,0.44,0.63}{{#1}}}
    \newcommand{\VerbatimStringTok}[1]{\textcolor[rgb]{0.25,0.44,0.63}{{#1}}}
    \newcommand{\SpecialStringTok}[1]{\textcolor[rgb]{0.73,0.40,0.53}{{#1}}}
    \newcommand{\ImportTok}[1]{{#1}}
    \newcommand{\DocumentationTok}[1]{\textcolor[rgb]{0.73,0.13,0.13}{\textit{{#1}}}}
    \newcommand{\AnnotationTok}[1]{\textcolor[rgb]{0.38,0.63,0.69}{\textbf{\textit{{#1}}}}}
    \newcommand{\CommentVarTok}[1]{\textcolor[rgb]{0.38,0.63,0.69}{\textbf{\textit{{#1}}}}}
    \newcommand{\VariableTok}[1]{\textcolor[rgb]{0.10,0.09,0.49}{{#1}}}
    \newcommand{\ControlFlowTok}[1]{\textcolor[rgb]{0.00,0.44,0.13}{\textbf{{#1}}}}
    \newcommand{\OperatorTok}[1]{\textcolor[rgb]{0.40,0.40,0.40}{{#1}}}
    \newcommand{\BuiltInTok}[1]{{#1}}
    \newcommand{\ExtensionTok}[1]{{#1}}
    \newcommand{\PreprocessorTok}[1]{\textcolor[rgb]{0.74,0.48,0.00}{{#1}}}
    \newcommand{\AttributeTok}[1]{\textcolor[rgb]{0.49,0.56,0.16}{{#1}}}
    \newcommand{\InformationTok}[1]{\textcolor[rgb]{0.38,0.63,0.69}{\textbf{\textit{{#1}}}}}
    \newcommand{\WarningTok}[1]{\textcolor[rgb]{0.38,0.63,0.69}{\textbf{\textit{{#1}}}}}
    \makeatletter
    \newsavebox\pandoc@box
    \newcommand*\pandocbounded[1]{%
      \sbox\pandoc@box{#1}%
      % scaling factors for width and height
      \Gscale@div\@tempa\textheight{\dimexpr\ht\pandoc@box+\dp\pandoc@box\relax}%
      \Gscale@div\@tempb\linewidth{\wd\pandoc@box}%
      % select the smaller of both
      \ifdim\@tempb\p@<\@tempa\p@
        \let\@tempa\@tempb
      \fi
      % scaling accordingly (\@tempa < 1)
      \ifdim\@tempa\p@<\p@
        \scalebox{\@tempa}{\usebox\pandoc@box}%
      % scaling not needed, use as it is
      \else
        \usebox{\pandoc@box}%
      \fi
    }
    \makeatother

    % Define a nice break command that doesn't care if a line doesn't already
    % exist.
    \def\br{\hspace*{\fill} \\* }
    % Math Jax compatibility definitions
    \def\gt{>}
    \def\lt{<}
    \let\Oldtex\TeX
    \let\Oldlatex\LaTeX
    \renewcommand{\TeX}{\textrm{\Oldtex}}
    \renewcommand{\LaTeX}{\textrm{\Oldlatex}}
    % Document parameters
    % Document title
    \title{analyzing\_parquet\_files\_v6}
    
    
    
    
    
    
    
% Pygments definitions
\makeatletter
\def\PY@reset{\let\PY@it=\relax \let\PY@bf=\relax%
    \let\PY@ul=\relax \let\PY@tc=\relax%
    \let\PY@bc=\relax \let\PY@ff=\relax}
\def\PY@tok#1{\csname PY@tok@#1\endcsname}
\def\PY@toks#1+{\ifx\relax#1\empty\else%
    \PY@tok{#1}\expandafter\PY@toks\fi}
\def\PY@do#1{\PY@bc{\PY@tc{\PY@ul{%
    \PY@it{\PY@bf{\PY@ff{#1}}}}}}}
\def\PY#1#2{\PY@reset\PY@toks#1+\relax+\PY@do{#2}}

\@namedef{PY@tok@w}{\def\PY@tc##1{\textcolor[rgb]{0.73,0.73,0.73}{##1}}}
\@namedef{PY@tok@c}{\let\PY@it=\textit\def\PY@tc##1{\textcolor[rgb]{0.24,0.48,0.48}{##1}}}
\@namedef{PY@tok@cp}{\def\PY@tc##1{\textcolor[rgb]{0.61,0.40,0.00}{##1}}}
\@namedef{PY@tok@k}{\let\PY@bf=\textbf\def\PY@tc##1{\textcolor[rgb]{0.00,0.50,0.00}{##1}}}
\@namedef{PY@tok@kp}{\def\PY@tc##1{\textcolor[rgb]{0.00,0.50,0.00}{##1}}}
\@namedef{PY@tok@kt}{\def\PY@tc##1{\textcolor[rgb]{0.69,0.00,0.25}{##1}}}
\@namedef{PY@tok@o}{\def\PY@tc##1{\textcolor[rgb]{0.40,0.40,0.40}{##1}}}
\@namedef{PY@tok@ow}{\let\PY@bf=\textbf\def\PY@tc##1{\textcolor[rgb]{0.67,0.13,1.00}{##1}}}
\@namedef{PY@tok@nb}{\def\PY@tc##1{\textcolor[rgb]{0.00,0.50,0.00}{##1}}}
\@namedef{PY@tok@nf}{\def\PY@tc##1{\textcolor[rgb]{0.00,0.00,1.00}{##1}}}
\@namedef{PY@tok@nc}{\let\PY@bf=\textbf\def\PY@tc##1{\textcolor[rgb]{0.00,0.00,1.00}{##1}}}
\@namedef{PY@tok@nn}{\let\PY@bf=\textbf\def\PY@tc##1{\textcolor[rgb]{0.00,0.00,1.00}{##1}}}
\@namedef{PY@tok@ne}{\let\PY@bf=\textbf\def\PY@tc##1{\textcolor[rgb]{0.80,0.25,0.22}{##1}}}
\@namedef{PY@tok@nv}{\def\PY@tc##1{\textcolor[rgb]{0.10,0.09,0.49}{##1}}}
\@namedef{PY@tok@no}{\def\PY@tc##1{\textcolor[rgb]{0.53,0.00,0.00}{##1}}}
\@namedef{PY@tok@nl}{\def\PY@tc##1{\textcolor[rgb]{0.46,0.46,0.00}{##1}}}
\@namedef{PY@tok@ni}{\let\PY@bf=\textbf\def\PY@tc##1{\textcolor[rgb]{0.44,0.44,0.44}{##1}}}
\@namedef{PY@tok@na}{\def\PY@tc##1{\textcolor[rgb]{0.41,0.47,0.13}{##1}}}
\@namedef{PY@tok@nt}{\let\PY@bf=\textbf\def\PY@tc##1{\textcolor[rgb]{0.00,0.50,0.00}{##1}}}
\@namedef{PY@tok@nd}{\def\PY@tc##1{\textcolor[rgb]{0.67,0.13,1.00}{##1}}}
\@namedef{PY@tok@s}{\def\PY@tc##1{\textcolor[rgb]{0.73,0.13,0.13}{##1}}}
\@namedef{PY@tok@sd}{\let\PY@it=\textit\def\PY@tc##1{\textcolor[rgb]{0.73,0.13,0.13}{##1}}}
\@namedef{PY@tok@si}{\let\PY@bf=\textbf\def\PY@tc##1{\textcolor[rgb]{0.64,0.35,0.47}{##1}}}
\@namedef{PY@tok@se}{\let\PY@bf=\textbf\def\PY@tc##1{\textcolor[rgb]{0.67,0.36,0.12}{##1}}}
\@namedef{PY@tok@sr}{\def\PY@tc##1{\textcolor[rgb]{0.64,0.35,0.47}{##1}}}
\@namedef{PY@tok@ss}{\def\PY@tc##1{\textcolor[rgb]{0.10,0.09,0.49}{##1}}}
\@namedef{PY@tok@sx}{\def\PY@tc##1{\textcolor[rgb]{0.00,0.50,0.00}{##1}}}
\@namedef{PY@tok@m}{\def\PY@tc##1{\textcolor[rgb]{0.40,0.40,0.40}{##1}}}
\@namedef{PY@tok@gh}{\let\PY@bf=\textbf\def\PY@tc##1{\textcolor[rgb]{0.00,0.00,0.50}{##1}}}
\@namedef{PY@tok@gu}{\let\PY@bf=\textbf\def\PY@tc##1{\textcolor[rgb]{0.50,0.00,0.50}{##1}}}
\@namedef{PY@tok@gd}{\def\PY@tc##1{\textcolor[rgb]{0.63,0.00,0.00}{##1}}}
\@namedef{PY@tok@gi}{\def\PY@tc##1{\textcolor[rgb]{0.00,0.52,0.00}{##1}}}
\@namedef{PY@tok@gr}{\def\PY@tc##1{\textcolor[rgb]{0.89,0.00,0.00}{##1}}}
\@namedef{PY@tok@ge}{\let\PY@it=\textit}
\@namedef{PY@tok@gs}{\let\PY@bf=\textbf}
\@namedef{PY@tok@ges}{\let\PY@bf=\textbf\let\PY@it=\textit}
\@namedef{PY@tok@gp}{\let\PY@bf=\textbf\def\PY@tc##1{\textcolor[rgb]{0.00,0.00,0.50}{##1}}}
\@namedef{PY@tok@go}{\def\PY@tc##1{\textcolor[rgb]{0.44,0.44,0.44}{##1}}}
\@namedef{PY@tok@gt}{\def\PY@tc##1{\textcolor[rgb]{0.00,0.27,0.87}{##1}}}
\@namedef{PY@tok@err}{\def\PY@bc##1{{\setlength{\fboxsep}{\string -\fboxrule}\fcolorbox[rgb]{1.00,0.00,0.00}{1,1,1}{\strut ##1}}}}
\@namedef{PY@tok@kc}{\let\PY@bf=\textbf\def\PY@tc##1{\textcolor[rgb]{0.00,0.50,0.00}{##1}}}
\@namedef{PY@tok@kd}{\let\PY@bf=\textbf\def\PY@tc##1{\textcolor[rgb]{0.00,0.50,0.00}{##1}}}
\@namedef{PY@tok@kn}{\let\PY@bf=\textbf\def\PY@tc##1{\textcolor[rgb]{0.00,0.50,0.00}{##1}}}
\@namedef{PY@tok@kr}{\let\PY@bf=\textbf\def\PY@tc##1{\textcolor[rgb]{0.00,0.50,0.00}{##1}}}
\@namedef{PY@tok@bp}{\def\PY@tc##1{\textcolor[rgb]{0.00,0.50,0.00}{##1}}}
\@namedef{PY@tok@fm}{\def\PY@tc##1{\textcolor[rgb]{0.00,0.00,1.00}{##1}}}
\@namedef{PY@tok@vc}{\def\PY@tc##1{\textcolor[rgb]{0.10,0.09,0.49}{##1}}}
\@namedef{PY@tok@vg}{\def\PY@tc##1{\textcolor[rgb]{0.10,0.09,0.49}{##1}}}
\@namedef{PY@tok@vi}{\def\PY@tc##1{\textcolor[rgb]{0.10,0.09,0.49}{##1}}}
\@namedef{PY@tok@vm}{\def\PY@tc##1{\textcolor[rgb]{0.10,0.09,0.49}{##1}}}
\@namedef{PY@tok@sa}{\def\PY@tc##1{\textcolor[rgb]{0.73,0.13,0.13}{##1}}}
\@namedef{PY@tok@sb}{\def\PY@tc##1{\textcolor[rgb]{0.73,0.13,0.13}{##1}}}
\@namedef{PY@tok@sc}{\def\PY@tc##1{\textcolor[rgb]{0.73,0.13,0.13}{##1}}}
\@namedef{PY@tok@dl}{\def\PY@tc##1{\textcolor[rgb]{0.73,0.13,0.13}{##1}}}
\@namedef{PY@tok@s2}{\def\PY@tc##1{\textcolor[rgb]{0.73,0.13,0.13}{##1}}}
\@namedef{PY@tok@sh}{\def\PY@tc##1{\textcolor[rgb]{0.73,0.13,0.13}{##1}}}
\@namedef{PY@tok@s1}{\def\PY@tc##1{\textcolor[rgb]{0.73,0.13,0.13}{##1}}}
\@namedef{PY@tok@mb}{\def\PY@tc##1{\textcolor[rgb]{0.40,0.40,0.40}{##1}}}
\@namedef{PY@tok@mf}{\def\PY@tc##1{\textcolor[rgb]{0.40,0.40,0.40}{##1}}}
\@namedef{PY@tok@mh}{\def\PY@tc##1{\textcolor[rgb]{0.40,0.40,0.40}{##1}}}
\@namedef{PY@tok@mi}{\def\PY@tc##1{\textcolor[rgb]{0.40,0.40,0.40}{##1}}}
\@namedef{PY@tok@il}{\def\PY@tc##1{\textcolor[rgb]{0.40,0.40,0.40}{##1}}}
\@namedef{PY@tok@mo}{\def\PY@tc##1{\textcolor[rgb]{0.40,0.40,0.40}{##1}}}
\@namedef{PY@tok@ch}{\let\PY@it=\textit\def\PY@tc##1{\textcolor[rgb]{0.24,0.48,0.48}{##1}}}
\@namedef{PY@tok@cm}{\let\PY@it=\textit\def\PY@tc##1{\textcolor[rgb]{0.24,0.48,0.48}{##1}}}
\@namedef{PY@tok@cpf}{\let\PY@it=\textit\def\PY@tc##1{\textcolor[rgb]{0.24,0.48,0.48}{##1}}}
\@namedef{PY@tok@c1}{\let\PY@it=\textit\def\PY@tc##1{\textcolor[rgb]{0.24,0.48,0.48}{##1}}}
\@namedef{PY@tok@cs}{\let\PY@it=\textit\def\PY@tc##1{\textcolor[rgb]{0.24,0.48,0.48}{##1}}}

\def\PYZbs{\char`\\}
\def\PYZus{\char`\_}
\def\PYZob{\char`\{}
\def\PYZcb{\char`\}}
\def\PYZca{\char`\^}
\def\PYZam{\char`\&}
\def\PYZlt{\char`\<}
\def\PYZgt{\char`\>}
\def\PYZsh{\char`\#}
\def\PYZpc{\char`\%}
\def\PYZdl{\char`\$}
\def\PYZhy{\char`\-}
\def\PYZsq{\char`\'}
\def\PYZdq{\char`\"}
\def\PYZti{\char`\~}
% for compatibility with earlier versions
\def\PYZat{@}
\def\PYZlb{[}
\def\PYZrb{]}
\makeatother


    % For linebreaks inside Verbatim environment from package fancyvrb.
    \makeatletter
        \newbox\Wrappedcontinuationbox
        \newbox\Wrappedvisiblespacebox
        \newcommand*\Wrappedvisiblespace {\textcolor{red}{\textvisiblespace}}
        \newcommand*\Wrappedcontinuationsymbol {\textcolor{red}{\llap{\tiny$\m@th\hookrightarrow$}}}
        \newcommand*\Wrappedcontinuationindent {3ex }
        \newcommand*\Wrappedafterbreak {\kern\Wrappedcontinuationindent\copy\Wrappedcontinuationbox}
        % Take advantage of the already applied Pygments mark-up to insert
        % potential linebreaks for TeX processing.
        %        {, <, #, %, $, ' and ": go to next line.
        %        _, }, ^, &, >, - and ~: stay at end of broken line.
        % Use of \textquotesingle for straight quote.
        \newcommand*\Wrappedbreaksatspecials {%
            \def\PYGZus{\discretionary{\char`\_}{\Wrappedafterbreak}{\char`\_}}%
            \def\PYGZob{\discretionary{}{\Wrappedafterbreak\char`\{}{\char`\{}}%
            \def\PYGZcb{\discretionary{\char`\}}{\Wrappedafterbreak}{\char`\}}}%
            \def\PYGZca{\discretionary{\char`\^}{\Wrappedafterbreak}{\char`\^}}%
            \def\PYGZam{\discretionary{\char`\&}{\Wrappedafterbreak}{\char`\&}}%
            \def\PYGZlt{\discretionary{}{\Wrappedafterbreak\char`\<}{\char`\<}}%
            \def\PYGZgt{\discretionary{\char`\>}{\Wrappedafterbreak}{\char`\>}}%
            \def\PYGZsh{\discretionary{}{\Wrappedafterbreak\char`\#}{\char`\#}}%
            \def\PYGZpc{\discretionary{}{\Wrappedafterbreak\char`\%}{\char`\%}}%
            \def\PYGZdl{\discretionary{}{\Wrappedafterbreak\char`\$}{\char`\$}}%
            \def\PYGZhy{\discretionary{\char`\-}{\Wrappedafterbreak}{\char`\-}}%
            \def\PYGZsq{\discretionary{}{\Wrappedafterbreak\textquotesingle}{\textquotesingle}}%
            \def\PYGZdq{\discretionary{}{\Wrappedafterbreak\char`\"}{\char`\"}}%
            \def\PYGZti{\discretionary{\char`\~}{\Wrappedafterbreak}{\char`\~}}%
        }
        % Some characters . , ; ? ! / are not pygmentized.
        % This macro makes them "active" and they will insert potential linebreaks
        \newcommand*\Wrappedbreaksatpunct {%
            \lccode`\~`\.\lowercase{\def~}{\discretionary{\hbox{\char`\.}}{\Wrappedafterbreak}{\hbox{\char`\.}}}%
            \lccode`\~`\,\lowercase{\def~}{\discretionary{\hbox{\char`\,}}{\Wrappedafterbreak}{\hbox{\char`\,}}}%
            \lccode`\~`\;\lowercase{\def~}{\discretionary{\hbox{\char`\;}}{\Wrappedafterbreak}{\hbox{\char`\;}}}%
            \lccode`\~`\:\lowercase{\def~}{\discretionary{\hbox{\char`\:}}{\Wrappedafterbreak}{\hbox{\char`\:}}}%
            \lccode`\~`\?\lowercase{\def~}{\discretionary{\hbox{\char`\?}}{\Wrappedafterbreak}{\hbox{\char`\?}}}%
            \lccode`\~`\!\lowercase{\def~}{\discretionary{\hbox{\char`\!}}{\Wrappedafterbreak}{\hbox{\char`\!}}}%
            \lccode`\~`\/\lowercase{\def~}{\discretionary{\hbox{\char`\/}}{\Wrappedafterbreak}{\hbox{\char`\/}}}%
            \catcode`\.\active
            \catcode`\,\active
            \catcode`\;\active
            \catcode`\:\active
            \catcode`\?\active
            \catcode`\!\active
            \catcode`\/\active
            \lccode`\~`\~
        }
    \makeatother

    \let\OriginalVerbatim=\Verbatim
    \makeatletter
    \renewcommand{\Verbatim}[1][1]{%
        %\parskip\z@skip
        \sbox\Wrappedcontinuationbox {\Wrappedcontinuationsymbol}%
        \sbox\Wrappedvisiblespacebox {\FV@SetupFont\Wrappedvisiblespace}%
        \def\FancyVerbFormatLine ##1{\hsize\linewidth
            \vtop{\raggedright\hyphenpenalty\z@\exhyphenpenalty\z@
                \doublehyphendemerits\z@\finalhyphendemerits\z@
                \strut ##1\strut}%
        }%
        % If the linebreak is at a space, the latter will be displayed as visible
        % space at end of first line, and a continuation symbol starts next line.
        % Stretch/shrink are however usually zero for typewriter font.
        \def\FV@Space {%
            \nobreak\hskip\z@ plus\fontdimen3\font minus\fontdimen4\font
            \discretionary{\copy\Wrappedvisiblespacebox}{\Wrappedafterbreak}
            {\kern\fontdimen2\font}%
        }%

        % Allow breaks at special characters using \PYG... macros.
        \Wrappedbreaksatspecials
        % Breaks at punctuation characters . , ; ? ! and / need catcode=\active
        \OriginalVerbatim[#1,codes*=\Wrappedbreaksatpunct]%
    }
    \makeatother

    % Exact colors from NB
    \definecolor{incolor}{HTML}{303F9F}
    \definecolor{outcolor}{HTML}{D84315}
    \definecolor{cellborder}{HTML}{CFCFCF}
    \definecolor{cellbackground}{HTML}{F7F7F7}

    % prompt
    \makeatletter
    \newcommand{\boxspacing}{\kern\kvtcb@left@rule\kern\kvtcb@boxsep}
    \makeatother
    \newcommand{\prompt}[4]{
        {\ttfamily\llap{{\color{#2}[#3]:\hspace{3pt}#4}}\vspace{-\baselineskip}}
    }
    

    
    % Prevent overflowing lines due to hard-to-break entities
    \sloppy
    % Setup hyperref package
    \hypersetup{
      breaklinks=true,  % so long urls are correctly broken across lines
      colorlinks=true,
      urlcolor=urlcolor,
      linkcolor=linkcolor,
      citecolor=citecolor,
      }
    % Slightly bigger margins than the latex defaults
    
    \geometry{verbose,tmargin=1in,bmargin=1in,lmargin=1in,rmargin=1in}
    
    

\begin{document}
    
    \maketitle
    
    

    
    \section{Did the bus actually stop? All-days analysis
(v6)}\label{did-the-bus-actually-stop-all-days-analysis-v6}

\textbf{Question.} For each scheduled bus call across all available
days: did the bus stop at the stop, or drive past?

\textbf{Data.} Every dated bronze folder containing
\texttt{stop\_visits} and \texttt{vehicle\_positions}. The latest
complete GTFS snapshot supplies the timetable; the real-time files
define the analysis window.

\textbf{What changed from v4}

\begin{itemize}
\tightlist
\item
  \texttt{vehicle\_positions.speed} is empty in the whole file, so v4's
  stopped/passed test could never fire: every event with a nearby ping
  ended up as \texttt{near\_no\_speed}. Speed is now derived from
  consecutive pings of the same vehicle.
\item
  v4 matched GPS to a stop by line + time only. That can pick up a
  \emph{different} bus on the same line, or one far away (likely where
  the ``\textgreater1 km nearest ping'' rows came from). Now each trip
  is first matched to one specific vehicle, and only that vehicle's
  track is used.
\item
  All available dated real-time files are analyzed together, with daily
  summaries. The latest complete GTFS snapshot is loaded once to avoid
  duplicating timetable rows. The side analyses (alerts, route-type
  checks, ferry checks, delays) are gone.
\end{itemize}

\textbf{What changed in v6} (new section 7b)

\begin{itemize}
\tightlist
\item
  Every call now gets evidence from the ping interval around the stop,
  measured along the route, so the \texttt{no\_ping\_at\_stop} blind
  spot is no longer empty.
\item
  \texttt{time\_lost\_s} / \texttt{p\_excess}: how much time the bus
  lost around the stop compared with ping intervals that contain no
  stop.
\item
  \texttt{ts\_offset\_enter\_s} / \texttt{ts\_offset\_leave\_s}: how the
  feed's \texttt{visit\_time} compares with when GPS says the bus
  entered / left the stop zone.
\end{itemize}

    \subsection{How to read the result}\label{how-to-read-the-result}

The ET feed cannot answer the question by itself: it copies one time
into both \texttt{actual\_arrival} and \texttt{actual\_departure} (so
there is no dwell time), and a ``recorded'' call only says the system
registered the bus there. GPS is the only evidence of standing still.

GPS pings arrive about every 30 s, and a bus that stands at a stop for
10--20 s is often between two pings. Speed is measured over those 30 s
steps, so even a real stop rarely shows speed 0: the step averages
driving and standing. That is why ``driving past'' is judged against the
bus's own usual speed nearby, not a fixed number. The result is a set of
evidence levels, not a yes/no:

\begin{longtable}[]{@{}
  >{\raggedright\arraybackslash}p{(\linewidth - 2\tabcolsep) * \real{0.5000}}
  >{\raggedright\arraybackslash}p{(\linewidth - 2\tabcolsep) * \real{0.5000}}@{}}
\toprule\noalign{}
\begin{minipage}[b]{\linewidth}\raggedright
\texttt{gps\_class}
\end{minipage} & \begin{minipage}[b]{\linewidth}\raggedright
Meaning
\end{minipage} \\
\midrule\noalign{}
\endhead
\bottomrule\noalign{}
\endlastfoot
\texttt{stopped\_confirmed} & 2+ pings within \texttt{NEAR\_M} of the
stop that look stopped: standing (\texttt{SLOW\_MS}) or well below the
bus's usual speed there (\texttt{SLOW\_RATIO}) \\
\texttt{likely\_stopped} & 1 such ping within \texttt{NEAR\_M} of the
stop \\
\texttt{likely\_passed} & ping(s) within \texttt{NEAR\_M}, all at or
above \texttt{FAST\_RATIO} of the bus's usual speed \\
\texttt{near\_unclear} & ping within \texttt{NEAR\_M}, speed in between
(or no usable speed) \\
\texttt{no\_ping\_at\_stop} & the bus track exists, but no ping landed
within \texttt{NEAR\_M}: cannot tell \\
\texttt{no\_gps} & bus matched, but there is a GPS hole at the time of
the stop \\
\texttt{no\_vehicle} & no bus could be matched to the trip \\
\end{longtable}

A bus driving past a stop is normal when nobody is waiting or has asked
to get off. ``Passed'' is not a service failure.

    \subsection{0. Setup}\label{setup}

    \begin{Verbatim}[commandchars=\\\{\}]
Analyzing 24 days (2026-10-20 through 2026-11-12); GTFS: 2026-11-12
    \end{Verbatim}

    \subsection{1. Timetable (GTFS)}\label{timetable-gtfs}

Only what the analysis needs: bus routes, their trips, stop coordinates,
and the ordered stop calls of each trip. \texttt{occ} numbers repeated
calls at the same stop within one trip (loops), so a call is identified
by trip + stop + occurrence.

    \begin{Verbatim}[commandchars=\\\{\}]
297 bus routes | 11,404 bus trips | 358,813 bus stop calls
calls timetabled as 'does not stop': 31
OK   stop calls whose stop\_id is not in gtfs\_stops: 0 of 358,813
OK   stop calls without coordinates: 0 of 358,813
    \end{Verbatim}

    \begin{Verbatim}[commandchars=\\\{\}, frame=single, framerule=2mm, rulecolor=\color{outerrorbackground}]
\textcolor{ansi-red-intense}{\textbf{The Kernel crashed while executing code in the current cell or a previous cell. 
}}
\textcolor{ansi-red-intense}{\textbf{Please review the code in the cell(s) to identify a possible cause of the failure. 
}}
\textcolor{ansi-red-intense}{\textbf{Click <a href='https://aka.ms/vscodeJupyterKernelCrash'>here</a> for more info. 
}}
\textcolor{ansi-red-intense}{\textbf{View Jupyter <a href='command:jupyter.viewOutput'>log</a> for further details.}}
    \end{Verbatim}

    \subsection{\texorpdfstring{2. Real-time stop calls
(\texttt{stop\_visits})}{2. Real-time stop calls (stop\_visits)}}\label{real-time-stop-calls-stop_visits}

The feed repeats every call in many snapshots as predictions change.
Each dated file is reduced before the next is loaded to keep memory
bounded. One call = \texttt{data\_frame\_ref} + \texttt{journey\_ref} +
\texttt{stop\_order}. Across snapshots I keep the newest record, but
take \texttt{actual\_*} from an older snapshot if the newest one lacks
it. Ferries are excluded and the rest is checked against the bus
allow-list from GTFS.

Only the columns the analysis uses are read (v4 read 21, most of them
unused here).

    \begin{Verbatim}[commandchars=\\\{\}]
328,121,410 snapshot rows -> 3,067,811 distinct calls
kept bus calls at known stops: 3,062,857 (removed 4,954)
visit\_status
observed       3037453
cancelled        14940
unconfirmed      10464

calls with both actual times where arrival == departure: 100.0\%  (so the feed
has no dwell time)
    \end{Verbatim}

    \subsection{3. Scheduled calls in the window, and what the feed says
about
each}\label{scheduled-calls-in-the-window-and-what-the-feed-says-about-each}

Active services come from \texttt{calendar} and
\texttt{calendar\_dates}. GTFS times can run past 24:00; the feed labels
a trip that starts after midnight with the clock date of its start, so
each scheduled call gets that same label (\texttt{service\_date}) before
it is joined to the feed. Only calls whose scheduled time falls inside
the window are kept, so no day is judged on data we do not have.

    \begin{Verbatim}[commandchars=\\\{\}]
scheduled bus calls inside the window: 462,953 on 16,170 trips
    \end{Verbatim}

    
    \begin{Verbatim}[commandchars=\\\{\}]
                   calls  share
final\_status                   
observed          452924  0.978
trip\_not\_in\_feed    6275  0.014
cancelled           2445  0.005
unconfirmed         1309  0.003
    \end{Verbatim}

    
    
    \begin{Verbatim}[commandchars=\\\{\}]
final\_status  cancelled  observed  trip\_not\_in\_feed  unconfirmed   calls
service\_date                                                            
2026-11-10          520    148613              2002          499  151634
2026-11-11          795    151722              2019          542  155078
2026-11-12         1109    151197              2254          250  154810
2026-11-13           21      1392                 0           18    1431
    \end{Verbatim}

    
    \begin{itemize}
\tightlist
\item
  \texttt{observed}: the feed has an actual time for the call.
\item
  \texttt{cancelled}: the feed marks the journey or call as cancelled.
\item
  \texttt{unconfirmed}: a feed record exists but has no actual time yet
  (mostly calls right at the end of the window).
\item
  \texttt{stop\_not\_in\_feed}, \texttt{trip\_not\_in\_feed}: nothing in
  the feed. This is a coverage gap, not proof that the bus did not run.
\end{itemize}

Only \texttt{observed} calls go on to the GPS test: the recorded time is
what lets us find the right bus.

    \subsection{\texorpdfstring{4. GPS positions
(\texttt{vehicle\_positions})}{4. GPS positions (vehicle\_positions)}}\label{gps-positions-vehicle_positions}

A position is repeated in every snapshot until the bus reports a new
one, so I keep one row per vehicle + \texttt{last\_updated}. Coordinates
are converted to metres on a local flat grid (accurate to about a metre
at these distances). The \texttt{speed} column is empty, so speed is
derived from consecutive pings of the same vehicle (metres moved ÷
seconds), both for the step into a ping and the step out of it.
\texttt{speed\_min} is the slower of the two: a bus standing at a stop
usually has at least one slow step next to the ping. \texttt{cruise\_ms}
is the bus's usual speed around that moment (the 75th percentile of the
neighbouring step speeds, i.e.~how fast it goes when it is not standing)
and \texttt{ratio\_min\ =\ speed\_min\ /\ cruise\_ms} says how much
slower than usual the bus was at that ping: about 1 when driving
normally, clearly lower when it stood or crawled.

    \begin{Verbatim}[commandchars=\\\{\}]
speed values present in the raw file: 0 of 10,493,828
10,493,828 raw rows across 24 days -> 7,552,289 distinct positions | 487
vehicles | 287 lines
positions from 2026-10-20 14:09:53.471000+00:00 to 2026-11-12 23:59:54+00:00
line\_ref found in GTFS route\_id: 99.8\%

seconds between positions of the same vehicle:
count    7543634.0
mean          47.4
std          483.0
min            0.1
10\%           29.0
50\%           30.0
90\%           31.0
max        79797.0

derived speed, m/s:
count    7458623.00
mean           6.84
std            7.68
min            0.00
10\%            0.35
50\%            6.11
90\%           13.19
max         1491.41

pings with a usable ratio to the bus's usual speed: 92.6\%

GPS coverage per hour (look for holes):
    \end{Verbatim}

    
    \begin{Verbatim}[commandchars=\\\{\}]
                           pings  vehicles
hour\_utc                                  
2026-10-20 14:00:00+00:00     18         9
2026-10-20 15:00:00+00:00  14439       237
2026-10-20 16:00:00+00:00  16714       194
2026-10-20 17:00:00+00:00  14912       158
2026-10-20 18:00:00+00:00  12158       147
{\ldots}                          {\ldots}       {\ldots}
2026-11-12 19:00:00+00:00  12090       147
2026-11-12 20:00:00+00:00  10300       116
2026-11-12 21:00:00+00:00   9886       101
2026-11-12 22:00:00+00:00   8993        96
2026-11-12 23:00:00+00:00   4462        75

[562 rows x 2 columns]
    \end{Verbatim}

    
    \subsection{5. Which bus ran which
trip?}\label{which-bus-ran-which-trip}

The GPS has no trip id, only \texttt{vehicle\_id} and
\texttt{line\_ref}. For every trip, each candidate vehicle on that line
is placed (by interpolating its track) at the recorded time of every
stop of the trip. The right bus is next to the stops at those times, so
it wins with a small median distance. Another bus on the same line (the
one before or after) is hundreds of metres off. A trip is matched only
if the best vehicle is within \texttt{ASSIGN\_MAX\_M} on at least
\texttt{ASSIGN\_MIN\_STOPS} stops. \texttt{runner\_up\_m} shows how
clear the choice was.

    \begin{Verbatim}[commandchars=\\\{\}]
observed calls to test: 452,924 on 15,588 trips
trips with a matched vehicle: 13,387 of 15,588 (85.9\%)

median distance of the matched bus to the trip's stops (m):
count    13387.0
mean        37.3
std          8.5
min         13.1
50\%         36.5
90\%         47.9
99\%         62.7
max         98.4
trips where the runner-up is within 2x of the winner (ambiguous): 27
    \end{Verbatim}

    \subsection{6. Stopped or not?}\label{stopped-or-not}

For each observed call, the matched bus is looked at in a window of
\texttt{WIN\_S} seconds around the recorded time. A ping within
\texttt{NEAR\_M} of the stop is ``at the stop''; it looks stopped if the
bus was standing (\texttt{speed\_min\ \textless{}=\ SLOW\_MS}) or at
most \texttt{SLOW\_RATIO} of its usual speed, and looks like driving
past if it was at least \texttt{FAST\_RATIO} of its usual speed. Classes
are explained at the top of the notebook.

    
    \begin{Verbatim}[commandchars=\\\{\}]
                    calls  share
gps\_class                       
stopped\_confirmed   29268  0.065
likely\_stopped     151579  0.335
likely\_passed       17622  0.039
near\_unclear        16363  0.036
no\_ping\_at\_stop    164084  0.362
no\_gps              12791  0.028
no\_vehicle          61217  0.135
    \end{Verbatim}

    
    
    \begin{Verbatim}[commandchars=\\\{\}]
gps\_class     stopped\_confirmed  likely\_stopped  likely\_passed  near\_unclear  no\_ping\_at\_stop  no\_gps  no\_vehicle   calls  likely\_passed\_share
service\_date                                                                                                                                  
2026-11-10                 6528           32176           3852          3538            35921    6245       60353  148613                0.026
2026-11-11                11337           59396           6866          6400            64133    3160         430  151722                0.045
2026-11-12                11326           59743           6819          6355            63225    3306         423  151197                0.045
2026-11-13                   77             264             85            70              805      80          11    1392                0.061
    \end{Verbatim}

    
    \subsection{7. Sanity checks and
results}\label{sanity-checks-and-results}

Read these before trusting the classes:

\begin{enumerate}
\def\labelenumi{\arabic{enumi}.}
\tightlist
\item
  How many trips got a bus, and how clearly (section 5 printout).
\item
  Whether \texttt{likely\_stopped} and \texttt{likely\_passed} both
  occur at plausible places and times. If one class is almost empty, the
  thresholds are off.
\item
  The speed ratio right next to stops (left chart): stopped and passed
  pings should form two groups, with the dashed lines in the gap between
  them.
\end{enumerate}

    \begin{Verbatim}[commandchars=\\\{\}]
calls with a ping within 30 m of the stop: 216,259 of 452,924 (47.7\%)

slowest speed (m/s) among the pings at the stop, for those calls:
count    216192.00
mean          3.87
std           2.86
min           0.00
10\%           0.51
25\%           2.05
50\%           3.42
75\%           5.10
90\%           7.37
max          88.92

same as a share of the bus's usual speed:
count    214125.00
mean          0.45
std           0.28
min           0.00
10\%           0.08
25\%           0.25
50\%           0.42
75\%           0.63
90\%           0.86
max          13.13
    \end{Verbatim}

    \begin{center}
    \adjustimage{max size={0.9\linewidth}{0.9\paperheight}}{analyzing_parquet_files_v6_files/analyzing_parquet_files_v6_20_1.png}
    \end{center}
    { \hspace*{\fill} \\}
    
    \begin{Verbatim}[commandchars=\\\{\}]
stops where buses most often drove past (at least 10 conclusive calls):
    \end{Verbatim}

    
    \begin{Verbatim}[commandchars=\\\{\}]
             stop\_id  calls  stopped  passed       stop\_name  conclusive  passed\_share
2057  NSR:Quay:46369     96        0      10  Nedre Øksnavad          10         1.000
2047  NSR:Quay:46359     90        0      14          Meland          14         1.000
2090  NSR:Quay:46415    140        0      13         Merland          13         1.000
2594  NSR:Quay:47094    122        1      10       Lutsiskog          11         0.909
2701  NSR:Quay:47220     96        1       9        Storhaug          10         0.900
2091  NSR:Quay:46416    126        2      17         Merland          19         0.895
3531  NSR:Quay:48320    125        2      13     Andakrossen          15         0.867
2700  NSR:Quay:47219     90        2      12        Storhaug          14         0.857
1759  NSR:Quay:46002    116        4      22        Vågsgata          26         0.846
1660  NSR:Quay:45878    102        2      11   Visteveien 25          13         0.846
3536  NSR:Quay:48328     90        2      11       Postvegen          13         0.846
2194  NSR:Quay:46558    221        6      32        Laberget          38         0.842
2486  NSR:Quay:46939    147        3      15   Nautanes nord          18         0.833
2492  NSR:Quay:46948    147        3      15    Nautanes sør          18         0.833
2597  NSR:Quay:47097    122        2       9     Lutsivatnet          11         0.818
    \end{Verbatim}

    
    
    \begin{Verbatim}[commandchars=\\\{\}]
    stop\_sequence                   stop\_name                          visit\_time vehicle\_id  min\_dist\_m  n\_near  slowest\_near\_ms  slowest\_near\_ratio          gps\_class
0               1               Jernbaneveien 2026-11-10 04:47:32.317000000+00:00       4074    4.168315       1         0.008822            0.001367     likely\_stopped
1               2     Stavanger politistasjon 2026-11-10 04:48:50.299000000+00:00       4074   30.232224       0              NaN                 NaN    no\_ping\_at\_stop
2               3                 Statens hus 2026-11-10 04:49:19.310000000+00:00       4074   73.586838       0              NaN                 NaN    no\_ping\_at\_stop
3               4               Vålandsbakken 2026-11-10 04:49:54.335000000+00:00       4074   16.170792       1         9.735159            1.104730      likely\_passed
4               5              Lindahlsbakken 2026-11-10 04:50:17.316000000+00:00       4074   72.017896       0              NaN                 NaN    no\_ping\_at\_stop
5               6                  Strømsbrua 2026-11-10 04:52:10.323000000+00:00       4074    6.540639       2         0.012494            0.001479  stopped\_confirmed
6               7           Hillevågstunnelen 2026-11-10 04:53:22.322000000+00:00       4074   25.761305       1         2.186779            0.321575     likely\_stopped
7               8                    Hillevåg 2026-11-10 04:54:28.333000000+00:00       4074    3.992692       2         0.589669            0.078352  stopped\_confirmed
8               9             Kvaleberg skole 2026-11-10 04:55:26.323000000+00:00       4074    6.102876       1         7.890624            0.967747      likely\_passed
9              10                    Sjøhagen 2026-11-10 04:56:00.327000000+00:00       4074   22.922277       1         6.197042            0.760037       near\_unclear
10             11                    Eikeberg 2026-11-10 04:57:28.325000000+00:00       4074    2.705034       1         2.718116            0.316657     likely\_stopped
11             12                     Mariero 2026-11-10 04:58:29.344000000+00:00       4074    4.292308       1         2.491586            0.290267     likely\_stopped
12             13                Lyngnesveien 2026-11-10 05:00:03.342000000+00:00       4074    8.352319       1         4.222723            0.468552     likely\_stopped
13             14               Hinnasvingene 2026-11-10 05:01:29.335000000+00:00       4074    5.301269       1         4.276948            0.455229     likely\_stopped
14             15                  Skråtårnet 2026-11-10 05:03:03.361000000+00:00       4074    8.788076       1         4.946945            0.526543     likely\_stopped
15             16               Stadionparken 2026-11-10 05:03:51.334000000+00:00       4074   58.920090       0              NaN                 NaN    no\_ping\_at\_stop
16             17                       Jåttå 2026-11-10 05:04:42.347000000+00:00       4074    2.251302       1         2.700244            0.280162     likely\_stopped
17             18              Gausel sentrum 2026-11-10 05:06:03.349000000+00:00       4074   47.160487       0              NaN                 NaN    no\_ping\_at\_stop
18             19                 Gauselhagen 2026-11-10 05:06:58.337000000+00:00       4074   20.406823       1        10.753728            1.056655      likely\_passed
19             20              Gausel stasjon 2026-11-10 05:08:12.356000000+00:00       4074  115.828239       0              NaN                 NaN    no\_ping\_at\_stop
20             21              Nådlandsbråtet 2026-11-10 05:09:14.346000000+00:00       4074   17.946101       1         5.112693            0.487576     likely\_stopped
21             22                   2020 park 2026-11-10 05:10:43.306000000+00:00       4074    1.903336       1         4.525619            0.471510     likely\_stopped
22             23                Løwenstrasse 2026-11-10 05:11:43.311000000+00:00       4074    4.888894       1         2.406803            0.243189     likely\_stopped
23             24         Jakob Askelands vei 2026-11-10 05:12:54.320000000+00:00       4074   14.354816       1         3.421828            0.339717     likely\_stopped
24             25                Porsmyrveien 2026-11-10 05:13:24.303000000+00:00       4074    9.029248       1        10.197917            1.012444      likely\_passed
25             26                     Kvadrat 2026-11-10 05:14:13.311000000+00:00       4074  140.526278       0              NaN                 NaN    no\_ping\_at\_stop
26             27                  Seljeveien 2026-11-10 05:14:51.308000000+00:00       4074   57.396127       0              NaN                 NaN    no\_ping\_at\_stop
27             28               Gamle Postvei 2026-11-10 05:16:07.312000000+00:00       4074    3.076219       1         5.658198            0.555898     likely\_stopped
28             29             Strandgata nord 2026-11-10 05:17:27.311000000+00:00       4074    4.860651       1         5.395898            0.513520     likely\_stopped
29             30                  Strandgata 2026-11-10 05:18:46.319000000+00:00       4074  108.799764       0              NaN                 NaN    no\_ping\_at\_stop
30             31               Jernbaneveien 2026-11-10 05:19:44.311000000+00:00       4074  125.351594       0              NaN                 NaN    no\_ping\_at\_stop
31             32               Olav Vs plass 2026-11-10 05:20:45.316000000+00:00       4074   26.883566       1         2.461860            0.233955     likely\_stopped
32             33  Sandnes bussterminal Ruten 2026-11-10 05:21:42.844645200+00:00       4074    7.104772       1         4.136262            0.581009             no\_gps
    \end{Verbatim}

    
    \subsection{7b. Evidence along the route
(v6)}\label{b.-evidence-along-the-route-v6}

v5 only counts a ping that lands within \texttt{NEAR\_M} of the stop, so
43\% of calls end up as \texttt{no\_ping\_at\_stop}. Here every call
gets evidence from the \textbf{ping interval that surrounds the stop},
measured \textbf{along the route} (GTFS shape, or the stop-to-stop line
if there is no shape):

\begin{enumerate}
\def\labelenumi{\arabic{enumi}.}
\tightlist
\item
  Each ping is snapped to the route, giving a distance along the route
  \texttt{s}. Jitter that makes \texttt{s} go backwards is removed
  (isotonic regression), because a bus does not drive backwards.
\item
  For each ping interval:
  \texttt{time\_lost\ =\ (time\ between\ the\ pings)\ -\ (distance\ covered\ /\ the\ bus\textquotesingle{}s\ usual\ speed)}.
  A bus that stood at a stop inside the interval loses time.
\item
  The same number is computed for ping intervals that contain \textbf{no
  stop} (traffic lights, turns, congestion stay in that comparison).
  \texttt{p\_excess} = share of stop-free intervals that lost at least
  as much time as this stop. Low = looks like a stop. High = no
  stop-like time loss, which is \textbf{not} proof of ``passed''.
\item
  Linear interpolation gives the time the bus enters and leaves the stop
  zone (\texttt{R\_STOP\_M} along the route). Compared with the feed's
  \texttt{visit\_time}, that tells whether the ET time is arrival or
  departure, and how accurate it is.
\end{enumerate}

\subsubsection{Comparing the feed with
GPS}\label{comparing-the-feed-with-gps}

The feed's \texttt{call\_type\ ==\ "recorded"} means ET marked the call
as recorded; it does not prove that the bus physically stopped. For a
side-by-side comparison, 7b marks a GPS call as \textbf{stop-like} when
\texttt{p\_excess\ \textless{}=\ 0.05}. This means its delay is rare
among stop-free intervals, not that there is a 95\% probability the bus
stopped. \textbf{ET only} means ET marked it recorded but GPS did not
find unusually high delay; \textbf{GPS only} means GPS found unusually
high delay but ET did not mark it recorded. These are disagreements
between indicators, not proof either method is wrong. The comparison
also shows how the older nearby-ping classes compare.
\texttt{likely\_passed} is not treated as the opposite of a recorded
call: a bus can pass a stop when nobody needs service.

I tested Newton and monotone-cubic interpolation against plain linear on
simulated data: no accuracy gain at 30 s pings, so only linear is used.
\texttt{dwell\_s} is not estimated: a 10-20 s stop between two pings
cannot be seen.

    \begin{Verbatim}[commandchars=\\\{\}]
route geometry: 5,386 trips from GTFS shapes, 0 from the stop-to-stop line, 0
without any
391,707 calls processed | 265,055 stop-free ping intervals for the comparison
    \end{Verbatim}

    \begin{center}
    \adjustimage{max size={0.9\linewidth}{0.9\paperheight}}{analyzing_parquet_files_v6_files/analyzing_parquet_files_v6_26_0.png}
    \end{center}
    { \hspace*{\fill} \\}
    
    \subsection{8. Results stay in memory}\label{results-stay-in-memory}

    \begin{Verbatim}[commandchars=\\\{\}]
Analysis complete. Results remain in notebook memory; no files were written.
    \end{Verbatim}

    \subsection{Limits}\label{limits}

\begin{itemize}
\item
  \textbf{30-second pings.} A short dwell can fall between two pings, so
  \texttt{no\_ping\_at\_stop} is a blind spot, not a ``passed''. Only
  \texttt{likely\_passed} is evidence of a bus driving past; absence of
  \texttt{stopped} is not.
\item
  \textbf{Standing near a stop is not the same as opening the doors.} A
  bus held at a red light within \texttt{NEAR\_M} of a stop looks the
  same as a bus that stopped. Door-open or passenger-count data would
  settle it; this dataset has neither.
\item
  \textbf{Matching depends on the recorded times.} The bus is found by
  where it was at the feed's actual times, so calls without an actual
  time (\texttt{unconfirmed}, \texttt{cancelled}, not in feed) cannot be
  tested.
\item
  \textbf{Thresholds are judgement.} \texttt{NEAR\_M},
  \texttt{SLOW\_MS}, \texttt{SLOW\_RATIO} and \texttt{FAST\_RATIO} are
  in the setup cell. They were set on simulated data, so check the
  histogram in section 7 on the real data and move the dashed lines into
  the gap between the two groups.
\item
  \textbf{\texttt{p\_excess} is a ranking score, not a strict test.}
  Traffic lights and congestion sit in the stop-free comparison, so it
  is conservative. ``No stop-like time loss'' is not proof that the bus
  drove past.
\item
  \textbf{The stop-free comparison needs enough stop-free intervals.}
  The notebook warns when there are fewer than \texttt{MIN\_NULL}. It
  also over-represents the stretches between stops.
\item
  \textbf{No dwell time.} A stop of 10-20 s between two ping intervals
  that look normal cannot be seen by any interpolation; only
  \texttt{time\_lost\_s} picks it up indirectly. \texttt{local\_gap\_s}
  shows the ping spacing behind each result.
\item
  \textbf{Route geometry.} Without GTFS shapes the route is the
  stop-to-stop straight line, which is coarser on curves (larger
  \texttt{MAX\_OFFROUTE\_STOPS}).
\end{itemize}


    % Add a bibliography block to the postdoc
    
    
    
\end{document}
